\documentclass[11pt]{article}
\usepackage[margin=1.1in]{geometry}
\usepackage{amsmath,amssymb,amsthm}
\usepackage[hidelinks]{hyperref}

\newcommand{\Lip}{\operatorname{Lip}}
\newcommand{\Free}{\mathcal F}
\newcommand{\pten}{\widehat{\otimes}_{\pi}}
\newtheorem{theorem}{Theorem}
\newtheorem{corollary}{Corollary}

\title{Literature Status Note for arXiv:1705.05145}
\author{fa\_banach\_001 / agent\_super\_00}
\date{2026-07-01}

\begin{document}
\maketitle

\section*{Status}

This is a lightweight literature-status packet, not a new proof.  It records
that two questions in Section 4 of Garcia-Lirola--Prochazka--Rueda Zoca,
\emph{A characterisation of the Daugavet property in spaces of Lipschitz
functions}, arXiv:1705.05145, were answered in later papers.

\section*{Source Questions}

The source paper asks in Section 4:
\[
  \text{if } M \text{ has property } (Z), \text{ must } M \text{ be a length space?}
\]
The same section ends with the vector-valued problem:
\[
  \Lip(M)\text{ Daugavet} \quad\Longrightarrow\quad
  \Lip(M,X)\text{ or } \Free(M)\pten X\text{ Daugavet?}
\]
for a pointed metric space \(M\) and a Banach space \(X\).

\section*{Supporting Answers}

Aviles--Martinez-Cervantes, \emph{Complete metric spaces with property (Z) are
length spaces}, arXiv:1808.00715, explicitly identifies the property-(Z)
problem as Question 1 from the source paper.  Its Main Theorem states:
\[
  M \text{ complete metric} \quad\Longrightarrow\quad
  M \text{ is length if and only if } M \text{ has property } (Z).
\]
Thus the complete metric-space form used in the source's Lipschitz-free and
Daugavet equivalences is already answered affirmatively.

Medina, \emph{A characterisation of the Daugavet property in spaces of
vector-valued Lipschitz functions}, arXiv:2305.05956, explicitly says that the
vector-valued problem was asked as Question 1 in the source paper and gives a
complete positive solution.  Theorem \(\ref{dummy-main}\) below records the
substance of Medina's Theorem \texttt{theo:mainDauga}, and Corollary
\(\ref{dummy-cor}\) records Medina's Corollary \texttt{cor:Daugavetvectorval}.

\begin{theorem}[Medina, arXiv:2305.05956, Theorem \texttt{theo:mainDauga}]
\label{dummy-main}
If \(M\) is a complete length metric space and \(X\) is a Banach space, then
\(\Lip(M,X)\) has the Daugavet property.
\end{theorem}

\begin{corollary}[Medina, arXiv:2305.05956, Corollary \texttt{cor:Daugavetvectorval}]
\label{dummy-cor}
If \(M\) is a length metric space and \(X\) is a nonzero Banach space, then
\((\Free(M)\pten X)^*\) has the Daugavet property; in particular,
\(\Free(M)\pten X\) has the Daugavet property.
\end{corollary}

Combining this with the source theorem that, for complete \(M\),
\(\Lip(M)\) has the Daugavet property if and only if \(M\) is length, answers
the vector-valued question for the complete-space setting of the source
characterization.

\section*{Scope}

The property-(Z) answer is recorded with the completeness hypothesis: this is
the form explicitly posed and answered by arXiv:1808.00715 and the form needed
for the source paper's Banach-space equivalences.  This packet does not assert
a blanket resolution of every non-complete formulation of ``length''.

The vector-valued answer is the special Lipschitz-free/vector-valued
Lipschitz result from arXiv:1705.05145.  It does not settle the broader
Werner problem about arbitrary projective tensor products of Daugavet spaces.

\begin{thebibliography}{9}

\bibitem{gpr18}
L. Garcia-Lirola, A. Prochazka, and A. Rueda Zoca,
\emph{A characterisation of the Daugavet property in spaces of Lipschitz
functions}, J. Math. Anal. Appl. 464 (2018), 473--492; arXiv:1705.05145.

\bibitem{am18}
A. Aviles and G. Martinez-Cervantes,
\emph{Complete metric spaces with property (Z) are length spaces},
arXiv:1808.00715.

\bibitem{medina23}
R. Medina,
\emph{A characterisation of the Daugavet property in spaces of vector-valued
Lipschitz functions}, arXiv:2305.05956.

\end{thebibliography}

\end{document}
